\documentclass{ctexbeamer}
\usecolortheme{spruce}
\setbeamertemplate{navigation symbols}{}
\setbeamertemplate{footline}[frame number]

\usepackage{booktabs}

\title[简短标题]{中文 Beamer 汇报模板}
\subtitle{用 \LaTeX{} 做学术报告}
\author{汇报人：张三}
\institute{示例大学 · 计算机学院}
\date{\today}

\begin{document}

\begin{frame}
  \titlepage
\end{frame}

\begin{frame}{目录}
  \tableofcontents
\end{frame}

\section{研究背景}

\begin{frame}{为什么用 Beamer 写中文幻灯片？}
  \begin{itemize}
    \item \texttt{ctexbeamer} 自动处理中文字体（Xe\LaTeX{} 引擎）
    \item 数学公式、代码、图表与论文完全一致
    \item 纯文本格式，Git 友好，随时改随时编
  \end{itemize}
  \begin{block}{一句话}
    改一处全局生效，不用再手动对齐文本框。
  \end{block}
\end{frame}

\section{方法与实验}

\begin{frame}{方法概览}
  \begin{columns}
    \begin{column}{0.5\textwidth}
      \begin{itemize}
        \item 输入：数据 $x \in \mathbb{R}^n$
        \item 模型：$f_\theta(x)$
        \item 目标：$\min_\theta \mathcal{L}(\theta)$
      \end{itemize}
    \end{column}
    \begin{column}{0.5\textwidth}
      \begin{table}
        \centering
        \begin{tabular}{lcc}
          \toprule
          方法 & 准确率 & F1 \\
          \midrule
          基线 & 81.2 & 0.79 \\
          本文 & \textbf{88.1} & \textbf{0.87} \\
          \bottomrule
        \end{tabular}
      \end{table}
    \end{column}
  \end{columns}
\end{frame}

\begin{frame}
  \centering
  {\LARGE 谢谢！}\\[1em]
  欢迎批评指正
\end{frame}

\end{document}
